% Сборка: lualatex -halt-on-error -interaction=nonstopmode -output-directory=SVD SVD/SVD_cur.tex
\documentclass[11pt,a4paper]{article}
\usepackage{fontspec}
\setmainfont{Noto Serif}
\usepackage[russian]{babel}
\usepackage[margin=20mm]{geometry}
\usepackage{amsmath,amssymb,microtype}
\setlength{\parindent}{0pt}
\setlength{\parskip}{3pt}
\setlength{\emergencystretch}{2em}
\newcommand{\R}{\mathbb R}
\newcommand{\Id}{\mathrm{Id}}
\newcommand{\T}{^{\mathsf T}}
\newcommand{\F}{_{\mathrm F}}
\DeclareMathOperator{\rank}{rank}
\begin{document}
\begin{center}
{\Large\bfseries Текущая реализация SVD-решателя}\\[4pt]
{\small По \texttt{ADP/solver/SVD.py}; обозначения \texttt{SVD\_solver.tex}}
\end{center}

\textbf{1. Вход и локальная оценка.}
Функция \texttt{solve} выполняет один шаг для многоиндексного базиса
$P\in\R^{m\times d}$ и CPU-статистик $U=(U_j)\in\R^{J\times p\times d}$,
$I=(I_j)\in\R^{J\times p}$. Проверяются входы; при отсутствии
\texttt{mass} используются единичные массы $N_j$, иначе массы должны быть
конечными и положительными. Базис ортонормируется. Через
\texttt{\_local\_refit} вычисляются minimum-norm решения
$g_j=\arg\min_g\|I_j-U_jP\T g\|_2^2$ (пакетное SVD,
с \texttt{lstsq} для чувствительных к потере ранга центров).
Эти $g_j$ затем фиксированы на весь внутренний шаг: полного чередования
локальных и глобальных оценок внутри \texttt{solve} нет.

\textbf{2. Два режима и начальное состояние.}
При $\lambda=\texttt{lambda\_prox}\ge0$ минимизируется
$F(B)=\sum_jN_j\|I_j-U_jB\T g_j\|_2^2+\lambda\|B-P\|\F^2$.
По умолчанию \texttt{low\_rank\_target="matrix"} ограничивает
$\rank(B)\le r$, $1\le r<m$. Режим \texttt{"correction"}
ограничивает $\rank(B-P)\le r$, $0\le r<m$.
Для единого описания положим
\[
 (X,I_j^0,P^0)=
 \begin{cases}
 (B,I_j,P),&\text{matrix},\\
 (\Delta,I_j-U_jP\T g_j,0),\quad B=P+\Delta,&\text{correction}.
 \end{cases}
\]
В обоих случаях начинают с $X=0$ и пустых факторов
$X=A\operatorname{diag}(\sigma)V\T$, $A\T A=V\T V=\Id_k$.
Кэш $W_j=U_jV$ позволяет считать предсказания без сборки $X$ на каждом шаге.
\texttt{\_FlatU} хранит непрерывный массив $U$ и его представление
$(Jp)\times d$ для матричных действий, при необходимости копируя вход.

\textbf{3. Инициализация очередной компоненты.}
Для $k=1,\ldots,r$ вычисляются
\[
 e_j=I_j^0-W_j\operatorname{diag}(\sigma)A\T g_j,\qquad
 Q=\sum_jN_jg_j(U_j\T e_j)\T+\lambda(P^0-X).
\]
Код строит $Q_\perp=Q-(QV)V\T$ и находит ведущий собственный вектор $u$
малой матрицы $Q_\perp Q_\perp\T$. При ведущем значении $\mu>0$
начальное направление равно $v=Q_\perp\T u/\sqrt\mu$;
при $\mu\le0$ следует остановка \texttt{zero\_gradient}.
Так новое начало выбирается вне уже найденного правого пространства.
Ортогональность к $V$ не навязывается последующим попеременным итерациям.

\textbf{4. Попеременное уточнение и контроль линейного решения.}
При фиксированном $Q$ повторяются нормированные решения
\[
 \left[\sum_jN_j\|U_jv\|_2^2g_jg_j\T+\lambda\Id_m\right]\widetilde a=Qv,
 \quad a=\widetilde a/\|\widetilde a\|_2,
\]
\[
 H_a\widetilde v=Q\T a,\qquad
 H_a=\sum_jN_j(a\T g_j)^2U_j\T U_j+\lambda\Id_d,
 \quad v=\widetilde v/\|\widetilde v\|_2.
\]
Малые системы решаются Холецким, с запасным \texttt{lstsq}.
При $\lambda>0$ и $d\le128$ по умолчанию система для $\widetilde v$
также решается Холецким; взвешенный дизайн собирается блоками до 16 MiB.
Иначе используется матрично-свободный LSMR для
$\min_x\|A_ax-t\|_2^2+\lambda\|x-z\|_2^2$, где
$(A_ax)_j=\sqrt{N_j}(a\T g_j)U_jx$,
$t_j=\sqrt{N_j}e_j$, $z=(P^0-X)\T a$.
\texttt{warm\_start=True} переиспользует предыдущий $\widetilde v$
внутри поиска пары; расширенная задача поправки сохраняет центр штрафа $z$.
Опция \texttt{precondition\_v} масштабирует оба блока по диагонали
$A_a\T A_a+\lambda\Id_d$; \texttt{adaptive\_krylov} меняет допуск LSMR.

Поиск заканчивается при $1-|v_{t+1}\T v_t|<\texttt{inner\_tol}$
(по умолчанию $10^{-6}$) либо после \texttt{inner\_maxiter} шагов (20).
Последнее ненормированное решение проверяется по
\[
 \eta=\frac{\|A_a\T(A_a\widetilde v-t)+\lambda(\widetilde v-z)\|_2}
 {\max(1,\|A_a\T t+\lambda z\|_2)}
 \le\max(10^{-7},10\,\texttt{inner\_tol}).
\]
При неудаче ускоренных вариантов выполняется строгое LSMR-уточнение;
оставшаяся неприемлемая невязка вызывает ошибку. Затем $a$ ещё раз точно
переоценивается для последнего $v$, даже если достигнут лимит итераций.

\textbf{5. Добавление, сжатие и принятие.}
Вычисляются $R=a\T Qv$, $D=\sum_jN_j(a\T g_j)^2\|U_jv\|_2^2+\lambda$,
масштаб $R/D$ и ожидаемый выигрыш $\gamma=R^2/D$.
При неположительном $D$ следует \texttt{zero\_curvature};
при $\gamma/\max(1,F)<\texttt{rank\_tol}$ (по умолчанию $10^{-8}$) ---
\texttt{rank\_tolerance}. Иначе к факторам добавляется компонента,
выполняются сокращённые QR и SVD малой матрицы произведения QR-факторов.
Если новое сингулярное значение не превышает
$\epsilon\max(m,d)s_{\max}$, следует \texttt{no\_rank\_growth}.
Кэш $U_jV$ преобразуется через QR-факторы; при
$\operatorname{cond}(R_R)>10^6$ пересчитывается непосредственно.

Для новых ортонормированных $a_\ell,v_\ell$ совместно решается $M\sigma=h$:
\begin{align*}
 M_{\ell t}&=\sum_jN_j(a_\ell\T g_j)(a_t\T g_j)
 (U_jv_\ell)\T(U_jv_t)+\lambda\delta_{\ell t},\\
 h_\ell&=\sum_jN_j(a_\ell\T g_j)(I_j^0)\T U_jv_\ell
                  +\lambda a_\ell\T P^0v_\ell.
\end{align*}
Кандидат принимается при конечном функционале и росте не более
$64\epsilon\max(1,F)$, где $\epsilon$ --- машинная точность; иначе возникает ошибка.
После принятия потеря численного ранга также завершает цикл.
История функционала, выигрыши, ранги, остановки и невязки сохраняются.

\textbf{6. Возврат базиса и коэффициентов.}
\texttt{solve\_fixed\_coefficients} возвращает сырую $B$ и диагностику.
\texttt{solve} в режиме matrix берёт активные правые направления по убыванию
$|\sigma_\ell|$ и дополняет их до $m$ из SVD матрицы
$P-(PV_{\rm act})V_{\rm act}\T$; при нулевом ранге возвращает $P$.
В режиме correction строится QR базиса $(P+\Delta)\T$ с проверкой полного
ранга строк; при нулевой поправке сохраняется $P$.
На итоговом базисе заново оцениваются локальные коэффициенты.
Возвращается \texttt{HPAOResult}: базис, коэффициенты, диагностика,
включая взвешенную ошибку данных и число локальных потерь ранга.
Убывание внутреннего $F$ не гарантирует его убывание после ортонормирования
или глобальную оптимальность; дополнение из $P$ не восстанавливает новые EDR-направления.

\textbf{Опциональная спектральная метрика.}
\texttt{metric\_power=0} оставляет описанный штраф Фробениуса.
При $\nu=\texttt{metric\_power}\in\{1/2,1\}$ код задаёт
$K=\alpha^2\Id_d+P\T\Lambda P$ (full) либо
$K=\alpha^2(\Id_d-P\T P)+P\T\Lambda P$ (orthogonal),
$\mathsf A=\rho\Id_d+(1-\rho)(K/\|K\|_2)^\nu$,
где $\rho=\texttt{metric\_floor}$. Требуются положительная определённость,
$\alpha>0$ и нормированный спектр $\Lambda$.
Штраф становится $\lambda\|(B-P)\mathsf A^{1/2}\|\F^2$;
тот же цикл применяется к $B\mathsf A^{1/2}$, $P\mathsf A^{1/2}$,
$U_j\mathsf A^{-1/2}$. Затем координаты возвращаются обратно,
и базис формируется в исходном пространстве. Плотная $d\times d$ метрика
не строится, но хранится дополнительный массив размера $U$.
\end{document}
